\chapter{Architectural Decomposition of \hydrot}
\label{chap-arch-decomp}

\phantomsection
\fancyhead[LE]{\thepage}
\fancyhead[RE]{\textbf{Hydrological Twin - Architectural Decomposition}}
\fancyhead[RO]{\thepage}
\fancyhead[LO]{\textbf{Hydrological Twin - Architectural Decomposition}}
\renewcommand{\headrulewidth}{1pt}

\setlength{\parindent}{0cm}
\setlength{\parskip}{9pt}

% ---------------------------------------------------------------------

\section{Global Architecture Overview}
\label{ch:architecture}

The Hydrological Twin architecture is structured as a hierarchy of git-synchronized scientific objects. Identity is no longer managed by an \texttt{IdCard}; it is defined through deterministic Git provenance and environment fingerprinting. Every object that participates in the twin space-state must expose a \texttt{GitSynchroRecord}. This transformation removes symbolic identity in favor of computational ontology.

\subsection{Architectural Principles}

The architecture is governed by five principles:

\begin{enumerate}
\item \textbf{Deterministic provenance.} All objects are reproducible through Git commit hashes and dependency locks.
\item \textbf{Layer symmetry.} Estimation, analysis and cartographic layers exist both at twin and compartment levels.
\item \textbf{Strict separation of concerns.} Model, Data and Numerical storage are distinct but synchronized.
\item \textbf{Append-only registry.} Validation events are immutable and hash-chained.
\item \textbf{Priority of reproducibility over convenience.}
\end{enumerate}

\subsection{Layered Architecture and Design Rationale}

The architecture decomposes into:

\begin{itemize}
\item \textbf{HydroTwin}
\item \textbf{Compartment}
\item \textbf{Simulation}
\item \textbf{Model Layer}
\item \textbf{Data Layer}
\item \textbf{NumPy Hypercube}
\item \textbf{Estimation Layer}
\item \textbf{Analysis Layer}
\item \textbf{Cartographic Layer}
\item \textbf{Universal Git Registry}
\end{itemize}

Each level introduces stricter reproducibility constraints.

\subsection{Inter-Layer Interactions and Data Flow}

The data flow is structured as:

\[
\text{Data} \rightarrow \text{Model} \rightarrow \text{Hypercube} \rightarrow \text{Estimation} \rightarrow \text{Analysis} \rightarrow \text{Cartography}
\]

Each transformation produces a new git-synchronized state.

\subsubsection{Managed Input-Output Flow}

\begin{figure}[H]
\centering
\fbox{%
\begin{minipage}{0.94\textwidth}
\centering
\textbf{Inputs}\\
public forcings, observations, spatial layers, model software, configuration manifests, private parameterizations, and entitlement proofs\\[0.2cm]
$\Downarrow$\\[0.2cm]
\textbf{Hydrological Twin managed services}\\
ingest $\rightarrow$ verify $\rightarrow$ structure $\rightarrow$ compute $\rightarrow$ analyse $\rightarrow$ render $\rightarrow$ register\\[0.2cm]
$\Downarrow$\\[0.2cm]
\textbf{Outputs}\\
public visualizations and aggregated indicators \quad | \quad private extracted hypercubes, simulation products, calibration results, and audit records
\end{minipage}}
\caption{Managed input-output flow of \hydrot.}
\label{fig:ht-managed-flow}
\end{figure}

\subsubsection{Semantic Contract of \hydrot}

\hydrot\ is not merely a file container or a software wrapper. Its semantics are those of a stateful scientific contract binding inputs, transformations, permissions, and outputs within the same auditable space-state. In this contract, observations, forcings, numerical models, parameters, and software provenance become synchronized components of one reproducible object rather than loosely coupled external resources.

This semantic contract is aligned with the current \texttt{HydrologicalTwinAlphaSeries} positioning: the standalone backend remains responsible for computation, data structuring, and controlled programmatic interfaces, whereas \cwv\ remains responsible for visualization, interaction, and GIS mediation. The White Book therefore treats \hydrot\ as the backend authority that determines what may be computed, exposed, exported, and registered.

An important semantic consequence is that NumPy hypercubes are internal numerical states, not public communication products. Public access is granted only to derived renderings, aggregated indicators, and non-reversible summaries that do not reveal raw hypercube values. All computational methods that would expose, transform, or regenerate raw hypercube content remain private services.

% --- MODIFICATION DE 2.1.3 ---
\subsubsection{Metadata Verification and Private-Key Authentication}
Before instantiation, HydrologicalTwin performs a multi-stage validation:
\begin{enumerate}
    \item \textbf{Metadata consistency check}: Verifies \texttt{GitSynchroRecord} hashes for code, data, and parameters.
    \item \textbf{Private-key authentication}: Required for operations beyond the \textit{Rendering} layer (e.g., \texttt{extract\_values}, simulation execution).
    \item \textbf{Immutability enforcement}: If metadata or key validation fails, the twin refuses to initialize and logs:
    \begin{verbatim}
    Error: Access denied. Invalid metadata or private-key.
    \end{verbatim}
\end{enumerate}

% --- FIN DE LA MODIFICATION ---



% ---------------------------------------------------------------------

\section{Access Semantics, Permission Tiers, and User Profiles}

\subsection{Permission tiers}

\hydrot\ implements a layered permission model in which the openness of visualization is preserved without exposing the computational core.

\begin{enumerate}
\item \textbf{Public rendering tier.} No authentication is required. Users may view maps, plots, dashboards, and aggregated statistics provided that the published products do not reveal raw NumPy hypercube values, full private parameterizations, or unrestricted extraction capabilities.

\item \textbf{Private computation tier.} Computational methods are private and accessible only after proof of purchase by cryptocurrency. The corresponding entitlement is tokenized on blockchain infrastructure and linked to a registered cryptographic credential used by the platform to authorize access. This tier unlocks raw extraction, advanced analysis, simulation execution, and controlled export services.

\item \textbf{Restricted expert tier.} Higher privacy levels can be attached to stronger tokenized entitlements for workflows such as parameter editing, self-fitting, calibration, Bayesian inference, multi-stage estimation, and administrative governance of private repositories.
\end{enumerate}

This tiering allows the same twin to support open scientific communication at the rendering level while preserving private economic and operational value at the computational level.

\subsection{Decorator-based enforcement strategy}

For an implementation that follows the White Book cases efficiently, the safest convention is a \textbf{fail-closed} one: the authenticated \texttt{HydroTwin} service class should be declared \emph{private by default}, and only the explicitly approved public methods should be marked as public.

\begin{itemize}
\item \textbf{Class decorator for the private computation tier.} A class-level decorator can declare that all methods of the authenticated service belong by default to the private computation tier and therefore require prior entitlement verification, user authentication, and audit logging.

\item \textbf{Method decorator for public exceptions.} Read-only methods intended for \cwv\ or other public rendering clients can then be whitelisted one by one through a method decorator. These methods should remain limited to renderings, aggregated indicators, masked values, and non-reversible summaries.

\item \textbf{Why this direction is preferable.} The inverse approach---public by default with private methods marked individually---is more fragile, because a forgotten decorator may expose a sensitive capability. In contrast, a private-by-default class with explicit public exceptions is aligned with the White Book distinction between a narrow public rendering surface and a broader private computation surface.
\end{itemize}

Decorators should nevertheless be understood as \textbf{policy declarations}, not as the security boundary itself. Effective protection must still be enforced at the service boundary through server-side authentication, authorization, entitlement verification, and systematic audit registration in the \texttt{GitSynchroRecord}-centred provenance chain.

In practice, \cwv\ should call only the methods exposed through a public service façade. Private clients should instead pass through an authenticated façade that validates credentials before any sensitive method is executed.

\subsection{User profiles}

The main user profiles are the following:

\begin{itemize}
\item \textbf{Public observer.} Consults public maps, plots, and summary indicators, with no access to raw hypercube data.
\item \textbf{Scientific analyst.} Uses token-gated private services to extract subsets, compare simulations and observations, and run reproducible analytical workflows.
\item \textbf{Operational operator.} Executes private simulation and scenario services for planning, monitoring, and operational decision support.
\item \textbf{Twin owner or commercial subscriber.} Holds the strongest entitlement level, including parameter management, calibration, and access to protected configuration assets.
\item \textbf{Auditor or regulator.} Reviews provenance records, validation traces, and published outputs without necessarily obtaining raw data modification rights.
\end{itemize}

These profiles are logical roles. A single actor may hold multiple roles if the corresponding entitlement tokens have been acquired and registered.

% ---------------------------------------------------------------------

\section{Core Object Hierarchy}

\subsection{HydroTwin}

\texttt{HydroTwin} is the top-level orchestrator.

It contains:

\begin{itemize}
\item \texttt{git\_synchro} of self
\item list of \texttt{Compartment}
\item twin-wide Estimation LayerOperatio
\item twin-wide Analysis Layer
\item twin-wide Cartographic Layer
\end{itemize}

% --- AJOUT DANS LA DESCRIPTION DE HydroTwin ---
HydroTwin enforces \textbf{role-based access control} via private-keys:
\begin{itemize}
    \item Acts as a gatekeeper for sensitive operations (e.g., data extraction, parameter modification).
    \item Logs all private actions with user-specific keys to ensure non-repudiation.
    \item Provides a \texttt{verify\_key()} method to validate user permissions dynamically.
\end{itemize}



% --- FIN DE L'AJOUT ---




\subsection{Compartment}

A \texttt{Compartment} represents a structured subspace of the hydrological system.

It contains:

\begin{itemize}
\item a list of \texttt{Simulation}
\item Estimation Layer
\item Analysis Layer
\item Cartographic Layer
\end{itemize}

\subsection{Simulation}

A \texttt{Simulation} binds:

\begin{itemize}
\item Model Layer (spatial support, temporal support, software git-synchro)
\item Data Layer (data manifest + parameter manifest)
\item NumPy Hypercube (with numpy version tracked)
\end{itemize}

% ---------------------------------------------------------------------

\section{Model Layer}
\label{ch:model-layer}

The Model Layer defines:

\begin{itemize}
\item Spatial support (CRS, resolution, mesh, topology)
\item Temporal support (time step, calendar, horizon)
\item Software provenance (Git commit, Python version, dependency lock)
\end{itemize}

The model is invalid if its git-synchro record cannot be resolved.

% ---------------------------------------------------------------------

\section{Data Layer}
\label{ch:data-layer}

The Data Layer includes:

\begin{itemize}
\item Forcings (public by default)
\item Model parameters (private by default)
\item Observational products
\end{itemize}

All data are content-addressed and hashed. Parameters are treated as data objects and therefore versioned equivalently.

% ---------------------------------------------------------------------

\section{Estimation Layer}
\label{ch:estimation-layer}

The Estimation Layer operates as a synchronization operator on the space-state.

It includes:

\begin{itemize}
\item Calibration
\item Bayesian inference
\item Data assimilation
\item Prior and likelihood definitions
\end{itemize}

Any estimation update produces a new git-synchronized twin state.

Priority rule: estimation can never modify data in place. It produces new states only.

% ---------------------------------------------------------------------

\section{Analysis Layer}
\label{ch:analysis-layer}

The Analysis Layer defines generic operators applicable across compartments.

Outputs must conform to a defined standard accepted by the twin.

\subsection{Transformation Operator}

Time transformations:

\begin{itemize}
\item Arithmetic mean over time
\item Monthly aggregation
\item Yearly aggregation
\end{itemize}

Spatial transformations:

\begin{itemize}
\item Spatial weighted arithmetic mean
\item Spatial harmonic mean (rates)
\item Spatial geometric mean (multiplicative processes)
\end{itemize}

\subsection{Extraction Operator}

The Extraction Operator produces NumPy hypercube subsets based on:

\begin{itemize}
\item spatial filters
\item temporal filters
\item variable selection
\end{itemize}

Extraction outputs are git-synchronized artifacts.

% ---------------------------------------------------------------------

\section{Cartographic Layer}
\label{ch:carto-layer}

The Cartographic Layer guarantees reproducible spatial rendering:

\begin{itemize}
\item CRS enforcement
\item Symbolization standards
\item Raster/vector export rules
\item QGIS compatibility
\end{itemize}

Map outputs must include provenance metadata and hash.

% ---------------------------------------------------------------------

\section{Git-Synchronized Registry}
\label{ch:registry}

\subsection{Universal Registry}

All validated \hydrot\ instances are appended to a hash-chained JSONL registry.

Each entry contains:

\begin{itemize}
\item UTC timestamp
\item git-synchro fingerprint
\item actor and host metadata
\item previous entry hash
\item entry hash
\end{itemize}

Tampering invalidates the chain.

\subsection{Priority of Actions in Registry Workflow}

Registry operations follow strict priority:

\begin{enumerate}
\item Validate twin consistency (all git-synchro records resolvable).
\item Freeze working tree (no dirty state allowed).
\item Append registry entry.
\item Commit registry file.
\item Push to remote.
\end{enumerate}

Failure at any stage aborts the process.

\subsection{Operational Rule}

% --- MODIFICATION DE 2.8.3 ---
A HydrologicalTwin instance becomes ontologically valid only after:
\begin{enumerate}
    \item Metadata verification (\texttt{verify\_metadata()}).
    \item \textbf{Private-key validation} (for restricted operations).
    \item Registry append (\texttt{append\_validated}).
    \item Git commit/push (\texttt{git\_commit}, \texttt{git\_push}).
\end{enumerate}
Public instances operate in a \textbf{read-only mode}, limited to the \textit{Rendering} layer.

Instantiating a Twin requires private-key clearance, as well as initiating a subTwin by extraction.

% --- FIN DE LA MODIFICATION ---



% ---------------------------------------------------------------------

\section{Action Prioritization Across the Architecture}

Execution priority across the entire system follows:

\begin{enumerate}
\item Git-synchro validation
\item Data manifest validation
\item Hypercube integrity check
\item Estimation execution
\item Analysis operators
\item Cartographic rendering
\item Registry append
\end{enumerate}

No analysis or cartographic output is authorized if provenance checks fail.

% ---------------------------------------------------------------------

\section{Architectural Consequence}

The architectural decomposition enforces:

\begin{itemize}
\item strict separation between model, data and estimation
\item deterministic reproducibility
\item explicit state transitions
\item cryptographic traceability
\end{itemize}

Hydrological Twin therefore operates not as a loose workflow, but as a formally structured git-synchronized space-state system in which every object, transformation and validation event is reproducible and audit-ready.
